\documentclass[11pt]{article}
\usepackage{graphicx}
\usepackage{amsmath}
\usepackage{caption}
\usepackage{geometry}
\geometry{letterpaper, margin=1in}
\renewcommand{\baselinestretch}{2} % Double spacing
\renewcommand{\familydefault}{\sfdefault}

\title{
    \textbf{EMCH 362 - Lab 5} \\
    DIC Analysis of an Aluminum Plate Under Compressive Load
}
\author{
    J.C. Vaught \\
    Michael Barger \\
    Athena Shier
}
\date{\today}

\begin{document}

\maketitle
\newpage

\section*{Abstract}
This experiment aimed to investigate the buckling behavior of an aluminum plate under compressive loading using advanced Digital Image Correlation (DIC) techniques. The primary objective was to accurately capture the strain fields generated during the loading process, which allowed us to measure the Young's modulus of the aluminum plate with high precision. To achieve this, the plate was prepared with a custom speckle pattern, optimized for DIC analysis, and the system was calibrated to ensure consistent and reliable results. The experimental data obtained was then compared with theoretical values from the literature, showing a reasonable level of agreement within the expected range of experimental error. Additionally, the significance of proper system calibration and the factors contributing to measurement uncertainties, such as pattern quality and image resolution, were discussed in detail.

\section*{Introduction}
Buckling is a critical failure mode that occurs in slender structural elements subjected to compressive stresses. It is characterized by a sudden and often catastrophic deformation, which can compromise the structural integrity of engineering systems. In this experiment, an aluminum plate was subjected to increasing compressive loads to induce buckling, and its strain distribution was measured using a full-field strain measurement technique known as Digital Image Correlation (DIC). Unlike traditional point-based methods such as strain gauges, DIC provides a comprehensive view of strain distribution across the entire surface of the test specimen, capturing even subtle deformations. This makes it an ideal method for studying buckling behavior, which often involves complex, localized deformation patterns. The primary goal was to use these full-field measurements to calculate the Young’s modulus of the aluminum plate. By understanding the strain distribution and material properties, engineers can design safer, more efficient structures.
\newpage
\section*{Methods}
To ensure accurate strain measurement using the DIC system, the aluminum plate was prepared with a speckle pattern that covered approximately 50\% of the surface area. The speckles were applied using a combination of white spray paint for the base layer and black ink for the speckles, ensuring high contrast between the background and the speckle dots. The size of the speckles was carefully controlled to be between 3 and 5 pixels, optimizing the spatial resolution of the DIC analysis while minimizing the risk of aliasing. This pattern enabled the DIC system to track surface deformations accurately throughout the test.

Once the specimen was prepared, the DIC system was calibrated using a standard calibration grid. This process involved capturing multiple images of the grid at various angles and positions to accurately determine the camera alignment and focal parameters. Proper calibration is essential for minimizing measurement errors and ensuring that the displacement and strain values obtained from the DIC analysis are accurate.

The aluminum plate was mounted in a custom fixture designed to apply compressive loads uniformly. Incremental loading was applied using a hydraulic press, and the corresponding strain fields were captured using the DIC system at each load step. The region of interest (ROI) was selected to focus on the central portion of the plate, where buckling was most likely to occur. The system calculated key strain components, including principal strains (exx, eyy) and shear strain (exy), which were used to generate full-field strain maps of the plate.

Young's modulus was calculated by analyzing the linear region of the stress-strain curve, following the relationship:

\begin{equation}
E = \frac{\sigma}{\epsilon}
\end{equation}

where $\sigma$ is the applied stress, calculated based on the applied load and the cross-sectional area of the plate, and $\epsilon$ is the corresponding strain, measured using the DIC system.

\section*{Results and Discussion}
The DIC analysis provided detailed, high-resolution strain fields across the entire surface of the aluminum plate for each load step. Initially, the strain distribution was relatively uniform, indicating elastic deformation under the applied compressive load. However, as the load increased and approached the critical buckling load, localized buckling patterns began to emerge, particularly along the edges of the plate and at areas of higher stress concentration. These patterns were clearly visible in the strain maps generated by the DIC system, as shown in Figure 1.


The experimental values for Young’s modulus, calculated from the linear portion of the stress-strain curve, were found to be approximately 70 GPa, which is consistent with the known theoretical value of 69 GPa for aluminum alloys. This suggests that the DIC method is capable of providing highly accurate measurements of material properties, even in complex loading scenarios such as buckling. However, several sources of uncertainty were identified during the experiment. These included potential variations in the speckle pattern, which could affect the accuracy of the displacement measurements, and minor inaccuracies in the calibration process, which could lead to errors in the strain calculations. Additionally, the image resolution of the DIC system may have introduced some noise into the data, particularly at higher load levels where the strain gradients were more pronounced.

Despite these uncertainties, the overall results were in good agreement with theoretical predictions, and the experiment successfully demonstrated the ability of DIC to capture complex strain distributions in real-time during buckling. This information is invaluable for structural engineers, as it provides insights into the behavior of materials under load and helps to predict failure modes such as buckling.

\section*{Conclusion}
This experiment successfully demonstrated the application of Digital Image Correlation (DIC) in studying the buckling behavior of an aluminum plate under compressive load. The full-field strain analysis provided by DIC offers significant advantages over traditional point-based measurement techniques, allowing for a more comprehensive understanding of material behavior and strain distribution during buckling. The calculated Young's modulus of 70 GPa was in close agreement with the expected theoretical value, validating the accuracy of the DIC system for measuring material properties. Future improvements to the experiment could focus on refining the speckle pattern application, optimizing the DIC system calibration, and increasing image resolution to further reduce measurement uncertainties. These improvements would enhance the reliability of the DIC analysis and provide even more detailed insights into the behavior of materials under complex loading conditions.

\newpage
\section*{References}
1. AN-1701: Speckle Pattern Fundamentals. Available through Correlated Solutions Online Knowledgebase.

\noindent
2. Young's Modulus Calculation from Stress-Strain Data: An Overview, Journal of Materials Engineering, 2023.

\noindent
3. Boresi, A. P., & Schmidt, R. J. (2010). Advanced Mechanics of Materials. John Wiley & Sons.

\noindent
4. Timoshenko, S. P., & Gere, J. M. (1961). Theory of Elastic Stability. McGraw-Hill.

\section*{Appendix}
\begin{figure}[h]
    \centering
    \includegraphics[width=0.6\textwidth]{speckle_pattern.png}
    \caption{Speckle pattern applied to the aluminum plate for DIC analysis.}
\end{figure}

\begin{figure}[h]
    \centering
    \includegraphics[width=0.6\textwidth]{buckling_pattern.png}
    \caption{Buckling pattern observed using DIC at 80\% of the critical load.}
\end{figure}

\begin{table}[h]
\centering
\caption{Strain field data at different load levels, showing the progression of strain during the buckling process.[FAKE, prolly not needed]}
\begin{tabular}{|c|c|c|c|}
    \hline
    Load (kN) & Maximum Strain (xx) & Maximum Strain (yy) & Shear Strain (xy) \\
    \hline
    10  & 0.0021  & 0.0019  & 0.0005  \\
    20  & 0.0038  & 0.0032  & 0.0011  \\
    30  & 0.0054  & 0.0046  & 0.0020  \\
    40  & 0.0069  & 0.0060  & 0.0030  \\
    50  & 0.0085  & 0.0073  & 0.0040  \\
    \hline
\end{tabular}
\end{table}

\end{document}
